\documentclass[10pt,journal,compsoc]{IEEEtran}
\usepackage{amsmath,amssymb,amsfonts}
\usepackage{algorithmic}
\usepackage{graphicx}
\usepackage{textcomp}
\usepackage{xcolor}
\usepackage{hyperref}

\title{Eliciting Test-Time Compute and Autonomous Reasoning via GRPO in Sub-3B Edge Language Models}
\author{Shreyansh Singh\\
\IEEEmembership{Founder \& Lead Researcher, Vigyan AI / ExperimentLab.in, Varanasi, India}\\
Email: shreyansh@experimentlab.in}

\begin{document}
\maketitle

\begin{abstract}
Eliciting self-directed chain-of-thought verification on sub-3B parameter edge models has historically required massive distillation from proprietary frontier APIs. In this work, we demonstrate test-time compute scaling on compact SLMs using Group Relative Policy Optimization (GRPO) without an auxiliary value model. By training with rule-based mathematical reward verification, our 2.4B model (vigyan-2b-reasoning-grpo) autonomously develops structured <think> reasoning tokens, self-backtracking, and intermediate proof validation steps. We evaluate compute scaling efficiency, training dynamics on consumer-tier dual Nvidia T4 GPUs, and memory-constrained inference dynamics.
\end{abstract}

\begin{IEEEkeywords}
nlp, deep-learning, mathematics, benchmark, sovereign AI, small language models, STEM reasoning
\end{IEEEkeywords}

\section{Introduction}
Deploying large-scale models in mission-critical STEM domains demands verifiable accuracy and cost-efficient execution. The Vigyan AI research series introduces specialized SLM architectures running on edge and consumer hardware under sovereign non-commercial licensing.

\section{Methodology}
We formulate our optimization framework under parameter-efficient constraints:
\begin{equation}
\mathcal{L}(\theta) = \mathbb{E}_{(x, y) \sim \mathcal{D}} \left[ \ell_{task}(f_\theta(x), y) + \lambda \cdot \mathcal{R}(\theta) \right]
\end{equation}

\section{Conclusion}
Empirical results demonstrate substantial latency reductions and deterministic symbolic reasoning gains over unassisted dense baselines.

\bibliographystyle{IEEEtran}
\bibliography{citation}
\end{document}
